\documentclass[11pt,a4paper]{article}
\usepackage[top=1.7cm,bottom=1.5cm,left=1.5cm,right=1.5cm,headsep=8pt,footskip=16pt]{geometry}
\usepackage{enumitem}
\usepackage{amssymb}
\usepackage{array}
\usepackage[table]{xcolor}
\usepackage{fancyhdr}
\usepackage{lastpage}
\usepackage{multicol}
\usepackage{polyglossia}
\setdefaultlanguage[calendar=gregorian, numerals=maghrib]{arabic}
\setotherlanguage{english}
\IfFontExistsTF{Amiri}{%
  \newfontfamily\arabicfont[Script=Arabic, Scale=1.08]{Amiri}
}{%
  \newfontfamily\arabicfont[Script=Arabic, Scale=1.08, Path=../fonts/, Extension=.ttf,
    UprightFont=*-Regular, BoldFont=*-Bold, ItalicFont=*-Italic,
    BoldItalicFont=*-BoldItalic]{Amiri}
}
\IfFontExistsTF{AlBattar}{%
  \newfontfamily\khatt[Script=Arabic, Scale=1.15]{AlBattar}
}{%
  \newfontfamily\khatt[Script=Arabic, Scale=1.15, Path=../fonts/,
    BoldFont=ae_AlBattar.ttf]{ae_AlBattar.ttf}
}
\renewcommand{\headrulewidth}{0.8pt}
\renewcommand{\footrulewidth}{0.8pt}
\fancyhf{}
\rhead{\khatt المدرسة العليا للأساتذة ببوزريعة}
\lhead{\khatt ملحقة بومرداس}
\lfoot{\khatt الأستاذ: د. بن عبد الرزاق يوسف}
\cfoot{\thepage/\pageref{LastPage}}
\pagestyle{fancy}
\parindent 0pt
\setlength{\columnsep}{0.6cm}
\setlength{\columnseprule}{0.4pt}
\setlist[enumerate,1]{label=\textbf{\arabic*.}, leftmargin=*, itemsep=5pt, topsep=2pt}
\begin{document}

\linespread{1.08}\selectfont
\begin{center}
{\LARGE\khatt\textbf{الاختبار الأول في الذكاء الاصطناعي}}\\[4pt]
{\normalsize وضعيات إشكالية ــ أسئلة متعددة الاختيارات ــ الفصول من 1 إلى 6}
\end{center}
\vspace{-2pt}\rule{\textwidth}{0.6pt}\\[2pt]
\textbf{التعليمة:} يتضمن الاختبار خمس وضعيات إشكالية. اقرأ كل وضعية بعناية، ثم اختر التحليل أو التصرف الأنسب
من بين الاختيارات الخمسة؛ لكل وضعية إجابة صحيحة واحدة فقط.
اكتب \textbf{حرف} الإجابة الصحيحة في جدول الإجابات في آخر الاختبار.
لكل وضعية 4 نقاط (المجموع 20 نقطة).
\vspace{4pt}

\noindent
{\large\khatt\textbf{الوضعية 1}}\hfill{\small(الفصل 1)}\par\nopagebreak\vspace{3pt}
\setlength{\fboxsep}{7pt}\colorbox{gray!10}{\parbox{\dimexpr\linewidth-14pt}{تُعِدّ أستاذة للسنة الرابعة متوسط درسًا في دراسة الأسلوب، وتحتاج إلى نص أصيل لكاتب جزائري معروف، مع أسئلة تحليلية تتدرج من الفهم إلى النقد. لديها محرك بحث، ومكتبة رقمية تضم مؤلفات الكاتب، وأداة ذكاء اصطناعي توليدي. وهي حريصة على الدقة العلمية. أي توزيع للأدوار بين هذه الوسائل هو الأنسب لطبيعة المهمة؟}}\par\nopagebreak
\begin{itemize}[leftmargin=1.0cm, labelsep=6pt, itemsep=3pt, topsep=4pt, parsep=0pt]
  \item[\textbf{أ)}] تطلب من الأداة التوليدية استرجاع نص الكاتب حرفيًا لأنها تدربت على كمّ هائل من المؤلفات المنشورة، ثم تطلب منها أسئلة متدرجة، وتكتفي بمقارنة سريعة لبداية النص بما يظهر في محرك البحث للتأكد من أنه النص نفسه.
  \item[\textbf{ب)}] تأخذ النص من المكتبة الرقمية لأن أصالته شرط في دراسة أسلوب الكاتب، ثم تستعين بالأداة التوليدية لاقتراح أسئلة متدرجة انطلاقًا من النص الذي تزوّدها به، وتراجع الأسئلة وتعدّلها بما يناسب أهداف الدرس ومستوى تلاميذها.
  \item[\textbf{ج)}] تختار النص من المكتبة الرقمية، ثم تكلّف الأداة التوليدية بتحديد الخصائص الأسلوبية للكاتب وتعتمد تحليلها مرجعًا للتصحيح، لأن الأداة اطلعت على الدراسات النقدية حول الكاتب، فيكون تحليلها أشمل من تحليل أي أستاذ بمفرده.
  \item[\textbf{د)}] تبحث في محرك البحث عن النص وتعتمد أول صفحة تنشره كاملًا، لأن المحرك يرتب النتائج حسب موثوقيتها، ثم تطلب من الأداة التوليدية تحليل أسلوبه وصياغة الأسئلة، دون الرجوع إلى المكتبة الرقمية التي تستغرق وقتًا أطول.
  \item[\textbf{هـ)}] تطلب من الأداة التوليدية كتابة نص جديد بأسلوب الكاتب نفسه، لأنه يحقق الهدف الأسلوبي ويتجنب مشكلات حقوق النشر، ثم تقدّمه للتلاميذ نموذجًا من أدب الكاتب، مع أسئلة تحليلية تولّدها الأداة وتراجعها الأستاذة بعناية.
\end{itemize}
\par\vspace{6pt}
\noindent
{\large\khatt\textbf{الوضعية 2}}\hfill{\small(الفصل 2)}\par\nopagebreak\vspace{3pt}
\setlength{\fboxsep}{7pt}\colorbox{gray!10}{\parbox{\dimexpr\linewidth-14pt}{درّب فريق نموذجًا لتصحيح أخطاء الهمزة على مئتي ألف جملة صحفية مصحَّحة، فبلغت دقته 97\% على جمل صحفية لم يرها أثناء التدريب، لكنها انخفضت إلى 61\% على إنشاءات تلاميذ الابتدائي. اقترح أعضاء الفريق تفسيرات وحلولًا مختلفة لهذه الفجوة. ويريد الفريق استعمال النموذج في المدارس قريبًا. أي اقتراح هو الأدق تفسيرًا والأنجع حلًّا؟}}\par\nopagebreak
\begin{itemize}[leftmargin=1.0cm, labelsep=6pt, itemsep=3pt, topsep=4pt, parsep=0pt]
  \item[\textbf{أ)}] الدقة المرتفعة على الجمل الصحفية تثبت أن النموذج تعلّم قواعد الهمزة تعلمًا عامًا، فالانخفاض راجع إلى رداءة خط التلاميذ عند الرقمنة؛ والحل تحسين جودة المسح الضوئي للإنشاءات، دون تغيير بيانات التدريب التي أثبتت فعاليتها.
  \item[\textbf{ب)}] النموذج حفظ الجمل الصحفية حفظًا آليًا ولم يتعلم شيئًا قابلًا للتعميم، بدليل انخفاض دقته؛ والحل تقليص بيانات التدريب إلى النصف لإجباره على التعميم، ثم إعادة قياس دقته على الجمل الصحفية نفسها التي تدرّب عليها.
  \item[\textbf{ج)}] الانخفاض طبيعي ولا يمكن تجاوزه لأن لغة الأطفال غير منتظمة؛ والحل الاكتفاء بالنموذج الحالي مع إعلام الأساتذة بدقته على الإنشاءات، لأن إضافة إنشاءات التلاميذ إلى التدريب ستُفسد ما تعلّمه النموذج من اللغة الصحفية السليمة.
  \item[\textbf{د)}] السبب اختلاف بيانات الاستعمال عن بيانات التدريب فعلًا، والحل إعادة تدريب النموذج على إنشاءات التلاميذ وحدها بعد حذف الجمل الصحفية كلها، ثم قياس دقته على الإنشاءات نفسها التي تدرّب عليها، لأنها أقرب عينة إلى الاستعمال الحقيقي.
  \item[\textbf{هـ)}] الفجوة سببها اختلاف بيانات الاستعمال عن بيانات التدريب، فأخطاء الأطفال وتراكيبهم نادرة في الصحافة؛ والحل إضافة عينة ممثلة من إنشاءات التلاميذ المصحَّحة إلى التدريب، ثم تقييم النموذج على إنشاءات أخرى لم يرها.
\end{itemize}
\par\vspace{6pt}
\noindent
{\large\khatt\textbf{الوضعية 3}}\hfill{\small(الفصل 3)}\par\nopagebreak\vspace{3pt}
\setlength{\fboxsep}{7pt}\colorbox{gray!10}{\parbox{\dimexpr\linewidth-14pt}{كتب أستاذ: «أنت أستاذ لغة عربية خبير. اكتب تمرينًا عن الحال»، فجاءت النتيجة عشرين جملة طويلة متفاوتة الصعوبة، بلا تصحيح، وبعضها يتضمن الحال الجملة وشبه الجملة، مع أن تلاميذه في السنة الثانية متوسط لم يدرسوا إلا الحال المفردة. ويريد طلبًا واحدًا يغنيه عن التعديل اليدوي. أي طلب معدَّل يعالج أسباب الخلل كلها؟}}\par\nopagebreak
\begin{itemize}[leftmargin=1.0cm, labelsep=6pt, itemsep=3pt, topsep=4pt, parsep=0pt]
  \item[\textbf{أ)}] «اكتب لتلاميذ السنة الثانية متوسط، الذين درسوا الحال المفردة فقط، عشرين جملة طويلة متنوعة الصعوبة يستخرجون منها الحال وصاحبها، ثم قدّم التصحيح في جدول، لأن كثرة الجمل وتفاوتها يغطيان كل الحالات الممكنة.»
  \item[\textbf{ب)}] «أنت أستاذ لغة عربية. اكتب تمرينًا من ست جمل قصيرة متقاربة الصعوبة عن الحال المفردة، يستخرج منها المتعلم الحال وصاحبها، ثم قدّم التصحيح في جدول، واجعل الجمل من الشعر الجاهلي لإثراء رصيد المتعلمين اللغوي.»
  \item[\textbf{ج)}] «أنت أستاذ لغة عربية خبير جدًا ودقيق. اكتب لتلاميذ السنة الثانية متوسط تمرينًا من ست جمل قصيرة متقاربة الصعوبة عن الحال بأنواعها، يستخرجون منها الحال وصاحبها، ثم قدّم التصحيح في جدول بعد التمرين.»
  \item[\textbf{د)}] «أنت أستاذ لغة عربية. اكتب لتلاميذ السنة الثانية متوسط، الذين درسوا الحال المفردة فقط، تمرينًا من ست جمل قصيرة متقاربة الصعوبة، يستخرجون منها الحال وصاحبها، ثم قدّم التصحيح في جدول بعد التمرين.»
  \item[\textbf{هـ)}] «أنت أستاذ لغة عربية. اكتب لتلاميذ السنة الثانية متوسط، الذين درسوا الحال المفردة فقط، تمرينًا قصيرًا يستخرجون فيه الحال وصاحبها من جمل مناسبة لمستواهم، واحرص على أن يكون التمرين واضحًا ومفيدًا وخاليًا من الأخطاء.»
\end{itemize}
\par\vspace{6pt}
\noindent
{\large\khatt\textbf{الوضعية 4}}\hfill{\small(الفصل 4)}\par\nopagebreak\vspace{3pt}
\setlength{\fboxsep}{7pt}\colorbox{gray!10}{\parbox{\dimexpr\linewidth-14pt}{سلّم طالب في السنة الثالثة بحثًا عن الموشحات الأندلسية، وأرفقه بفقرة يصرّح فيها بأنه استعمل أداة ذكاء اصطناعي «في التحرير»، لكن الأستاذ اكتشف أن ثلاثة مراجع من سبعة غير موجودة، وأن فقرات كاملة منقولة من الأداة دون تعديل. وطلب الطالب مقابلة لتوضيح موقفه. كيف يُقيَّم الموقف في ضوء مبادئ المسؤولية والشفافية والإفصاح؟}}\par\nopagebreak
\begin{itemize}[leftmargin=1.0cm, labelsep=6pt, itemsep=3pt, topsep=4pt, parsep=0pt]
  \item[\textbf{أ)}] الطالب أخلّ بمبدأ المسؤولية فقط بسبب المراجع المختلقة، أما الفقرات المنقولة دون تعديل فلا إشكال فيها مادام قد أفصح، لأن الإفصاح العام عن الاستعمال يغطي أي حجم من النص المولَّد في البحوث الجامعية.
  \item[\textbf{ب)}] يُلغى البحث ويُمنع الطالب من استعمال أي أداة مستقبلًا، لأن اجتماع الإفصاح الناقص والمراجع المختلقة يدل على نية الغش، والنية وحدها كافية للحكم دون الرجوع إلى نظام المؤسسة أو محاورة الطالب.
  \item[\textbf{ج)}] الإفصاح موجود لكنه غير دقيق، لأن عبارة «في التحرير» تخفي حجم الاستعمال فتُخلّ بالشفافية؛ والمراجع المختلقة يتحمل الطالب مسؤوليتها كاملة لأن التحقق منها واجبه، فيُعالَج الأمر وفق نظام المؤسسة مع محاورته وإعادة العمل.
  \item[\textbf{د)}] الطالب خالف مبدأ الإفصاح وحده لأنه لم يذكر اسم الأداة ونسختها، أما المراجع المختلقة فخطأ تقني تتحمله الشركة المطوِّرة، فيكفي أن يضيف اسم الأداة ونسختها ليصبح العمل مطابقًا لمبادئ الاستعمال الأخلاقي.
  \item[\textbf{هـ)}] الطالب احترم المبادئ الثلاثة لأنه أفصح عن الاستعمال، والإفصاح يُسقط عنه المسؤولية عن أخطاء الأداة؛ فيكفي أن يحذف المراجع المختلقة، ثم يُقيَّم البحث على أساس ما بقي من مراجع صحيحة وفقرات سليمة، كما تنص على ذلك أغلب الأنظمة الجامعية.
\end{itemize}
\par\vspace{6pt}
\noindent
{\large\khatt\textbf{الوضعية 5}}\hfill{\small(الفصل 5)}\par\nopagebreak\vspace{3pt}
\setlength{\fboxsep}{7pt}\colorbox{gray!10}{\parbox{\dimexpr\linewidth-14pt}{لخّصت أداة لأستاذة مقالًا من عشر صفحات عن عادات القراءة عند المراهقين في ثلاث فقرات. عند المقارنة وجدت ثلاث ملاحظات: نسبة مئوية غير واردة في المقال، وحذف فقرة عن القراءة الرقمية، وتحويل عبارة الكاتب «قد يرتبط بـ» إلى «يؤدي إلى». وتريد توزيعه على زملائها في الاجتماع القادم. أي تقييم للملخص هو الأدق؟}}\par\nopagebreak
\begin{itemize}[leftmargin=1.0cm, labelsep=6pt, itemsep=3pt, topsep=4pt, parsep=0pt]
  \item[\textbf{أ)}] الملخص غير أمين في موضعين على الأقل: إضافة نسبة لا أصل لها، وتحويل الارتباط إلى سببية يغيّر حكم الكاتب؛ أما حذف فقرة القراءة الرقمية فيُحكم عليه بحسب وزنها في المقال، فإن كانت فكرة رئيسية وجب إدراجها.
  \item[\textbf{ب)}] الملخص أمين في مجمله لأن أفكاره الرئيسية موجودة في المقال، والملاحظات الثلاث هفوات شكلية، فيكفي حذف النسبة وإضافة جملة عن القراءة الرقمية، دون مراجعة صيغ الأحكام لأنها لا تغيّر مضمون الملخص.
  \item[\textbf{ج)}] الخلل الأخطر هو حذف فقرة القراءة الرقمية، لأن الملخص يجب أن يغطي كل فقرات المقال بالتساوي، أما النسبة المضافة فتُقبل إن كانت معقولة، وتغيير صيغة العبارة تحسين أسلوبي لا يمسّ المعنى، فالقارئ يبحث عن الفكرة لا عن العبارة.
  \item[\textbf{د)}] الخلل الوحيد هو النسبة المضافة لأنها معلومة جديدة؛ أما تحويل «قد يرتبط» إلى «يؤدي إلى» فتصرف مقبول في تلخيص يقتضي الاختصار والوضوح، وحذف أي فقرة جائز دائمًا لأن الملخص بطبيعته أقصر من الأصل.
  \item[\textbf{هـ)}] الملخص غير أمين في موضعين: النسبة المضافة وحذف فقرة القراءة الرقمية؛ أما تحويل الارتباط إلى سببية فاستنتاج منطقي يحق للملخِّص الوصول إليه، لأن الارتباط القوي بين ظاهرتين يدل غالبًا على علاقة سببية بينهما، كما تقرر في مناهج التلخيص.
\end{itemize}
\par\vspace{6pt}
\begin{center}
{\khatt\textbf{جدول الإجابات}}\\[3pt]
\renewcommand{\arraystretch}{1.6}
\begin{tabular}{|>{\centering\arraybackslash}m{1.9cm}|*{5}{>{\centering\arraybackslash}m{1.2cm}|}}
\hline
\textbf{رقم السؤال} & \textbf{1} & \textbf{2} & \textbf{3} & \textbf{4} & \textbf{5} \\ \hline
\textbf{الإجابة} &  &  &  &  &  \\[8pt] \hline
\end{tabular}

\vspace{8pt}
\renewcommand{\arraystretch}{2.2}
\begin{tabular}{|>{\centering\arraybackslash}m{7.5cm}|>{\centering\arraybackslash}m{3.2cm}|>{\centering\arraybackslash}m{3.2cm}|}
\hline
\textbf{الاسم واللقب} & \textbf{الفوج} & \textbf{العلامة} \\ \hline
 & & \hfill /20 \\ \hline
\end{tabular}
\end{center}
\end{document}
